\documentclass[11pt]{article}
\usepackage[a4paper,margin=1in]{geometry}
\usepackage{amsmath,amssymb,amsthm,microtype}
\newtheorem{theorem}{Theorem}[section]
\theoremstyle{remark}\newtheorem{remark}[theorem]{Remark}
\title{Hardy-Dominant Annular Coefficient Conversion and Its Sharp Gap Order}
\author{Bin Wang}\date{July 2026}
\begin{document}\maketitle
\begin{abstract}
The direct $H^\infty$ coefficient estimate of RH-300 used Cauchy's bound and
has gap constant of order $\eta^{-1}$ when $\rho=Re^\eta$.  Since normalized
$H^\infty$ embeds contractively into $H^2$, the Hardy estimate improves this
to order $\eta^{-1/2}$.  We prove the Hardy constant exactly and show that
its square-root order is optimal: equality is attained in $H^2$, while
dyadic Rudin--Shapiro blocks give the same order on the $H^\infty$ unit ball.
At $R=1.4$, $\rho=1.41$, the constants are $139.007\ldots$ and
$8.292\ldots$.  No norm decay for the actual mismatch is asserted.
\end{abstract}

\section{Coefficient functional}
Let $f(z)=\sum_{n\ge2}b_nz^n$ be analytic on $|z|<\rho$, put
$x=R/\rho<1$, and define
\[
 P_R(f)=\sum_{n\ge2}|b_n|R^n.
\]
Boundary $H^2$ norms use normalized angular measure.

\section{Hardy-dominant theorem}
\begin{theorem}\label{thm:hardy}
For every $f\in H^2(\rho)$,
\[
 P_R(f)\le
 \|f\|_{H^2(\rho)}\frac{x^2}{\sqrt{1-x^2}}.
\]
In particular, for $f\in H^\infty(\rho)$,
\[
 P_R(f)\le
 \|f\|_{H^\infty(\rho)}\frac{x^2}{\sqrt{1-x^2}}.
\]
The $H^2$ constant is exact.  If $\rho=Re^\eta$, then the best uniform
order on the $H^\infty$ unit ball is also $\Theta(\eta^{-1/2})$ as
$\eta\downarrow0$.
\end{theorem}
\begin{proof}
Write $u_n=|b_n|\rho^n$.  Cauchy--Schwarz gives
\[
 P_R(f)=\sum_{n\ge2}u_nx^n
 \le\left(\sum_{n\ge2}u_n^2\right)^{1/2}
 \left(\sum_{n\ge2}x^{2n}\right)^{1/2},
\]
which is the first estimate.  The second uses
$\|f\|_{H^2}\le\|f\|_{H^\infty}$.  Equality in the Hilbert-space estimate
is obtained by taking the boundary-scaled coefficient vector proportional
to $(x^n)_{n\ge2}$.

For the $H^\infty$ lower order, let $N=2^m$ and choose Rudin--Shapiro signs
$\varepsilon_0,\ldots,\varepsilon_{N-1}$ with
\[
 \sup_{|w|=1}\left|\sum_{j=0}^{N-1}\varepsilon_jw^j\right|
 \le\sqrt{2N}.
\]
Then
\[
 f_N(z)=\frac1{\sqrt{2N}}
 \sum_{j=0}^{N-1}\varepsilon_j(z/\rho)^{j+2}
\]
has $H^\infty(\rho)$ norm at most one and
\[
 P_R(f_N)=
 \frac{x^2(1-x^N)}{\sqrt{2N}(1-x)}.
\]
Choose a dyadic $N$ with $(2\eta)^{-1}\le N\le\eta^{-1}$.  Since
$x=e^{-\eta}$, the last expression is bounded below by
$c\eta^{-1/2}$.  The Hardy upper bound is asymptotic to
$(2\eta)^{-1/2}$, proving the order.
\end{proof}

\section{Constants and protocol}
At $R=1.4$, $\rho=1.41$,
\[
 \frac{x^2}{1-x}=139.0070921986\ldots,\qquad
 \frac{x^2}{\sqrt{1-x^2}}=8.2924678943\ldots.
\]
The computation evaluates the Hardy scale and admissible dyadic
Rudin--Shapiro lower examples at $\eta=10^{-1},10^{-2},10^{-3}$, using
lengths $8,64,512$.  Non-dyadic lengths are rejected by the implementation
because the displayed normalization is then not certified by the standard
identity.

\begin{remark}
This theorem optimizes the conversion of a known norm into an absolute
coefficient budget.  It does not prove that the actual $g_\sigma$ norm
decays.  Gates A--E remain false/open and no Hilbert--Polya or RH statement
follows.
\end{remark}
\end{document}
